\documentclass[11pt]{amsart}
\usepackage[T1]{fontenc}
\usepackage[utf8]{inputenc}
\usepackage[margin=1in]{geometry}
\usepackage{amssymb,booktabs,array,longtable}
\usepackage[hidelinks]{hyperref}
\newtheorem{proposition}{Proposition}[section]
\DeclareMathOperator{\Spec}{Spec}
\DeclareMathOperator{\deh}{deh}
\newcommand{\OO}{\mathcal O}
\newcommand{\code}[1]{\texttt{#1}}
\title[Comparing proofs of Zariski's Main Theorem]{Finite section lifting and Zariski's Main Theorem:\\
a comparison with EGA, Horrocks, and the Peskine route}
\date{1 November 2026}
\begin{document}
\emergencystretch=1em
\begin{abstract}
The argument in \code{zmt-proof.tex} gives a constructive proof of the
finite-fibre algebraic form of Zariski's Main Theorem. Its conclusion
was already proved constructively by Alonso, Coquand and Lombardi.
Its cohomological organisation belongs to the formal-functions tradition,
but its finite section correction and extraction differ from the
Stein-factorisation proof in EGA and from the conductor argument used in
Stacks. We identify the precise classical antecedents and the adaptations
made for the checked proof. This comparison does not establish historical
priority for the particular combination.
\end{abstract}
\maketitle

\section{The statement and the two proof routes}

Let $R$ be a commutative ring, $I$ an ideal, and $E=R[x_1,\ldots,x_n]$
an algebra. The statement under discussion is
\[
 E/IE\text{ finite over }R/I
 \quad\Longrightarrow\quad
 \exists s\in1+IE:\quad
 B=R[s,sx_1,\ldots,sx_n]\text{ finite over }R,\quad B_s\simeq E_s.
\]
Here ``finite'' means finite as a module. There is no noetherian or
finite-presentation assumption. This is the conclusion of
Alonso--Coquand--Lombardi (ACL), Theorem~1.3 and Corollary~1.4
\cite{ACL}; their remark also removes injectivity of $R\to E$.
Thus neither this statement nor its constructive validity is new.

\paragraph{Peskine, ACL, and Stacks.}
ACL follows Peskine, using a conductor, strong transcendence and
subresultants to replace the minimal-prime argument, then induction on
algebra generators (\cite{ACL}, \S\S2--3, especially
Propositions~2.16--2.18 and Theorem~3.1).
Stacks \S10.123 likewise develops the conductor and strong-transcendence
route; its introduction explicitly cites Peskine, Raynaud and notes by
Coquand \cite{StacksZMT}. Our proof does not follow that chain of lemmas.
Peskine's original 1966 article was not inspected for this comparison:
the detailed comparison is with ACL and the actual Stacks proof.

\paragraph{EGA and the projective route.}
EGA~III$_1$, \S4.3, obtains Stein factorisation from formal functions and
coherence of proper direct images. Proposition~4.4.1 identifies the
isolated-fibre locus in the proper case; Theorem~4.4.3 applies this after
a projective compactification \cite{EGA}. This is a different route from the one
used in Stacks \S10.123. The relevant chain is
\[
 \text{formal functions}\ \longrightarrow\
 \text{Stein factorisation and connected fibres}\ \longrightarrow\
 \text{ZMT}.
\]
Our argument retains projective charts and section lifting, and replaces
this ending by explicit integrality certificates.

\section{The exact classical approximation principle}

Let $W$ be proper over a noetherian ring $A$, and $\mathcal F$ coherent.
For an ideal $I$ and a fixed $p$, Stacks Lemma~30.20.4(3) gives $c$ such that
\begin{align*}
 &\operatorname{im}\bigl(H^p(W,\mathcal F/I^k\mathcal F)
       \longrightarrow H^p(W,\mathcal F/I^r\mathcal F)\bigr)\\
 &\hspace{12mm}=\operatorname{im}\bigl(H^p(W,\mathcal F)
       \longrightarrow H^p(W,\mathcal F/I^r\mathcal F)\bigr)
       \qquad(k\ge r+c).
\end{align*}
This statement precedes the completion theorem \cite{StacksApprox}.
Its proof uses finiteness over the Rees algebra of graded cohomology and
of its obstruction submodule \cite{StacksFormal}.

This finite control is already explicit in EGA~III$_1$:
(4.1.7.12) makes the graded obstruction module finite,
(4.1.7.20) kills its sufficiently distant transition maps, and
(4.1.7.21) identifies the stable image with the image of genuine
cohomology. The comparison proof is credited there to Serre
(footnote on p.~126) \cite{EGA}. Consequently, replacing completion by a finite
approximation step is not itself a new idea.

For the application below, only $p=0$ and $r=1$ are needed.
One needs an actual section with the prescribed reduction modulo $I$;
one does not need an actual section inducing a prescribed entire
compatible formal section. Formal functions identifies a completion,
and does not assert that every element of that completion comes from
the uncompleted module.

\paragraph{The precise connection with Horrocks.}
In Theorem~1, pp.~714--715, Horrocks extends a bundle to $\mathbf P^1_A$
for a noetherian local ring $A$. After twisting, the special-fibre
splitting has nonnegative degrees and a trivial summand. Vanishing of
$H^1$ gives surjectivity through the infinitesimal thickenings; EGA's
formal-functions theorem then lifts a chosen special-fibre section to
a genuine section. That section splits off a summand and enables
induction on rank \cite{Horrocks}.
The common mechanism is lifting one useful residual section.
Horrocks does not give a ZMT proof here, and our ZMT proof does not
assume the analogous $H^1$ vanishing.

\section{A direct mathematical bridge to the certificate}

The following explains the comparison mathematically. It is our
reconstruction from the approximation principle, rather than an
attribution of this exact proof to EGA or Horrocks.

First suppose $R=A$ is noetherian. Put $X=\Spec E$ and choose a
projective compactification $X\subset W\subset\mathbf P^n_A$, with
$X=D_+(Z)$. Write $W_k=W\times_A A/I^k$ and
$X_k=X\times_A A/I^k$.

The algebra $E/I^kE$ is finite over $A/I^k$: choose a monic polynomial
$P_i$ for each $x_i$ modulo $IE$; then $P_i(x_i)^k=0$ modulo $I^kE$.
The open immersion $X_k\to W_k$ is also proper, because $X_k$ is
finite over $A/I^k$ and $W_k$ is separated. Thus $X_k$ is open and
closed in $W_k$. Extend the section $Z^{-1}$ of $\OO(-1)$ on $X_k$
by zero to its complement. These sections have the same reduction on
$W_1$.

Apply the approximation principle to $\mathcal F=\OO_W(-1)$.
It gives $\zeta\in H^0(W,\OO_W(-1))$ whose restriction to $X_1$ is
$Z^{-1}$. Therefore
\[
 s=(Z\zeta)|_X\in E,\qquad s\equiv1\pmod{IE}.
\]
On each chart $X_i\ne0$, represent $\zeta$ by a homogeneous fraction
of degree $-1$. Comparing it with the fraction on the chart $Z$ and
clearing denominators gives
\begin{equation}\label{DL}
 s x_i^{a_i+1}=q_i(x),\qquad \deg q_i\le a_i.
\end{equation}
The exponents may be increased to clear localisation equalities as
well as the displayed denominators. No injection of a chart ring into
a further localisation is required. Choose also a polynomial
expression $s=\sigma(x)$ with $\deg\sigma\le D$.

\begin{proposition}
The data \eqref{DL} imply that $s,sx_1,\ldots,sx_n$ are integral over $A$.
\end{proposition}
\begin{proof}
Let $F_d$ be the span of monomials of total degree at most $d$, and put
$d=D+\sum_i a_i$. For a monomial $x^\beta$ of degree at most $d+1$,
if $\beta_i>a_i$ for some $i$, use \eqref{DL} to express $sx^\beta$
with degree at most $|\beta|-1\le d$. Otherwise use $s=\sigma(x)$,
obtaining degree at most $D+\sum_i a_i=d$. Hence $sF_{d+1}\subseteq F_d$.
The finite module $V=F_d$ contains $1$ and is stable under both $s$
and every $sx_i$. Apply the determinant trick separately to these
endomorphisms of the same module $V$.
\end{proof}

Monomial reduction in the resulting monic equations makes
$B=A[s,sx_1,\ldots,sx_n]$ finite. Every $x_i$ is $(sx_i)/s$ in $E_s$,
and localisation preserves the injection $B\subseteq E$; hence
$B_s\simeq E_s$.

This ending explains the purpose of the negative twist: a denominator
of exponent $a_i+1$ has a numerator of degree $a_i$, giving exactly
the one-degree drop in \eqref{DL}. It needs one finite stable module,
without constructing the finite algebra $H^0(W,\OO_W)$ or proving
connectedness of fibres.

\section{What the constructive proof changes}

The checked proof develops the preceding mechanism using finite
polynomial and Laurent data. It first extracts finitely many residual
monic equations from $E/IE$. Their coefficients are specialised from
a polynomial ring $A_0=\mathbb Z[C_1,\ldots,C_M]$. This auxiliary ring
has the proved constructive coherence, discreteness and noetherian
properties needed internally. Transport of the final finite identities
then gives the result over arbitrary $R$, without deciding equality
or ideal membership in $R$.

\paragraph{Sections without decomposing schemes.}
Write $J$ for the homogeneous relations and $I=(b_1,\ldots,b_t)$.
The proof constructs cochains $z_r$ by explicit geometric-series
identities, with
\[
 dz_r\in I^{r+1}C^1+JC^1,\qquad
 z_r-z_0\in IC^0+JC^0,\qquad (z_r)_Z=Z^{-1}.
\]
These are the algebraic replacement for the zero-extended sections on
$W_{r+1}$. The formulas do not require chart idempotents or a choice
of a compatible infinite family of representatives.

\paragraph{One dependence relation instead of a uniform cohomology theorem.}
Over $D=A_0[U_1,\ldots,U_t]$, the evaluated Rees presentation uses
the actual kernel of
\[
 D[X,Z]\longrightarrow (A_0[X,Z]/J)[T],\qquad U_j\longmapsto b_jT.
\]
Its finite bihomogeneous generators and their $J$-certificates are
constructed. Image powers of $I$ are used, so no flatness assumption
or identification of $I^r\otimes M$ with $I^rM$ is made.

The obstruction cochains are raised through a bounded partial free
resolution. Monomial contractions leave only finitely many top
Laurent coefficients, in a finite free module $D^v$. Each obstruction
merely has such a partner. Noetherian induction obtains a dependence
by using one partner at a time inside the existential conclusion.
There is no choice of a whole sequence of partners.

After projection to one coefficient degree and evaluation $U_j=b_j$,
this gives, for some $N$, coefficients $r_i\in I^{N-i}$ and
$v\in I^{N+1}C^0$ such that
\[
 d\left(z_N-\sum_{i<N}r_i z_i-v\right)\in JC^1.
\]
The corrected section still agrees with $z_0$ modulo $I+J$, since
$N-i\ge1$. This includes $N=0$. Extraction gives \eqref{DL}.

Finite syzygy kernels and the Rees kernel are constructed in this
argument. What is avoided is a separate computation of generators
for the whole cohomology obstruction module, and a bound valid for
all approximate sections. The theorem asks only for one dependence
in the sequence at hand.

\paragraph{Assessment of originality.}
The ingredients---projective compactification, Rees finiteness,
\v Cech contractions, dimension shifting, denominator clearing and
the determinant trick---are established. The specific organisation
around one sequence, partners under mere existence, and a single
finite stable module is the feature worth documenting. The sources
examined here do not present that combination. That observation
supports describing the work as an \emph{alternative constructive
proof by finite section lifting}; it does not establish that the
combination has never appeared elsewhere.

For perspective, Zariski's 1949 analytical proof uses valuations,
completions of local domains and analytical irreducibility, notably
in Theorems~3--4 \cite{Zariski}. It is a more distant comparison than
the explicit section-lifting mechanism above. No exhaustive review
of Zariski's other papers, EGA~IV, or later textbooks was made here.

\section{Connection with the Unit development}

The mathematical proof is in \code{unit/doc/zmt-proof.tex}; the
current endpoint and scope are recorded in
\code{unit/sag/zmt-general/README.txt}. The correspondence is:

{\small
\begin{longtable}{>{\raggedright\arraybackslash}p{.35\textwidth}
                  >{\raggedright\arraybackslash}p{.59\textwidth}}
\toprule
Mathematical step & Directories below \code{unit/sag/}\\
\midrule
\endhead
Finite residual data and arbitrary-base descent & \code{zmt-descent}\\
Explicit approximate sections & \code{zmt-sections-general}\\
Finite dependence and weighted projection & \code{zmt-certificate-general},
 \code{zmt-weighted-general}\\
Rees membership and correction & \code{zmt-rees-general},
 \code{zmt-correction-general}\\
Degree-lowering relations & \code{zmt-extract-general}\\
One stable module and determinant trick & \code{zmt-certificate-extract-general}\\
All finite numbers of generators & \code{zmt-general/FinalZMTGeneral.seg}\\
\bottomrule
\end{longtable}}

The public result is
\begin{verbatim}
FinalZMTGeneral.Theorem[R,setR,I,propI,Ea,n,x,gen,fo].result
\end{verbatim}
It includes a finite spanning family for the actual image subalgebra
$B\subseteq E$ and an equivalence given by the canonical localisation
map. The noetherian auxiliary base and the dependence theorem are
internal, not hypotheses of this endpoint. The development's ambient
univalence, replacement and surjection parameters are unchanged;
it introduces neither excluded middle nor a choice principle.

This comparison reads the proof and its recorded formalisation status;
it does not report a new checker run. Its scope is the displayed
finite-fibre algebraic theorem. A global quasi-finite scheme
factorisation requires additional covering and gluing statements.
The classical reconstruction in Section~3 is explanatory; it is not
a separate Unit formalisation of schemes, Stein factorisation or
formal functions.

\begin{thebibliography}{9}
\bibitem{ACL}
M.~E. Alonso, T.~Coquand and H.~Lombardi,
\emph{Revisiting Zariski Main Theorem from a constructive point of view},
J.~Algebra \textbf{406} (2014), 46--68.
\url{https://arxiv.org/abs/1602.00597}.
The local \code{ACL.pdf} was read, especially Theorem~1.3,
Corollary~1.4 and \S\S2--3.

\bibitem{EGA}
A.~Grothendieck, with J.~Dieudonn\'e,
\emph{\'El\'ements de g\'eom\'etrie alg\'ebrique III, premi\`ere partie},
Publ. Math. IH\'ES \textbf{11} (1961), 5--167.
\url{https://www.numdam.org/articles/10.1007/BF02684273/}.
Original scan inspected at (4.1.7.12)--(4.1.7.21), pp.~127--129;
(4.2.1), pp.~129--130; (4.3.1)--(4.3.3), pp.~130--131;
(4.4.1)--(4.4.3), pp.~135--136. These are the volume page numbers,
not the alternative continuous pagination printed at the foot.

\bibitem{Horrocks}
G.~Horrocks, \emph{Projective modules over an extension of a local ring},
Proc. London Math. Soc. (3) \textbf{14} (1964), 714--718.
\url{https://doi.org/10.1112/plms/s3-14.4.714}.
Theorem~1 and its proof were read in the local \code{phorrocks.pdf}.

\bibitem{StacksZMT}
The Stacks Project, \S10.123, \emph{Zariski's Main Theorem}, Tag~00PI.
\url{https://stacks.math.columbia.edu/tag/00PI}.

\bibitem{StacksFormal}
The Stacks Project, \S30.20, \emph{The theorem on formal functions},
especially Lemmas~30.20.1 and 30.20.3.
\url{https://stacks.math.columbia.edu/tag/02O7}.

\bibitem{StacksApprox}
The Stacks Project, Lemma~30.20.4(3), Tag~02OB.
\url{https://stacks.math.columbia.edu/tag/02OB}.

\bibitem{Zariski}
O.~Zariski, \emph{A simple analytical proof of a fundamental property
of birational transformations}, Proc. Nat. Acad. Sci. USA
\textbf{35} (1949), 62--66.
\url{https://doi.org/10.1073/pnas.35.1.62}.
The paper was inspected through the NCBI scan/text.
\end{thebibliography}
\end{document}
